\documentclass[11pt]{amsart}

\usepackage[a4paper,margin=1.05in]{geometry}
\usepackage{amsmath,amssymb,amsthm}
\usepackage{stmaryrd}      % \llbracket \rrbracket
\usepackage{mathtools}
\usepackage{enumitem}
\usepackage{booktabs}
\usepackage{xcolor}
\usepackage[colorlinks=true,linkcolor=blue!45!black,citecolor=blue!45!black,urlcolor=blue!45!black]{hyperref}

% ---- notation ----------------------------------------------------------
\newcommand{\ctr}{\mathbin{\triangleright}}      % container: shapes > positions
\newcommand{\ext}[1]{\llbracket #1 \rrbracket}    % extension
\newcommand{\comp}{\mathbin{\triangleleft}}       % substitution / sequential composition
\newcommand{\Set}{\mathsf{Set}}
\newcommand{\Cont}{\mathsf{Cont}}
\newcommand{\DCont}{\mathsf{DCont}}
\newcommand{\Cat}{\mathsf{Cat}}
\newcommand{\Poly}{\mathsf{Poly}}
\newcommand{\fwd}{\mathsf{fwd}}
\newcommand{\bwd}{\mathsf{bwd}}
\newcommand{\bk}[1]{{#1}^{\sharp}}                % backward part of a morphism
\newcommand{\y}{\mathsf{y}}
\newcommand{\zs}{\bowtie}                          % Zappa--Szep weld
\DeclareMathOperator{\Store}{Store}
\DeclareMathOperator{\getop}{get}
\DeclareMathOperator{\putop}{put}
\DeclareMathOperator{\cofree}{Cofree}

\theoremstyle{plain}
\newtheorem{theorem}{Theorem}
\newtheorem{proposition}{Proposition}
\newtheorem{lemma}{Lemma}
\newtheorem{corollary}{Corollary}
\theoremstyle{definition}
\newtheorem{definition}{Definition}
\newtheorem{example}{Example}
\theoremstyle{remark}
\newtheorem{remark}{Remark}

\title[Tool-calling protocols are directed containers]{Tool-calling protocols are directed containers}
\author{MacBeth}
\address{Kodamai}
\email{}
\date{19 September 2026}

\begin{document}

\begin{abstract}
An agent tool-calling protocol -- the discipline by which a language-model agent issues a typed
request and receives a typed response -- is a container: its shapes are the request schemas and its
positions are the admissible responses. We observe that the \emph{harness} running such a protocol
is not a separate meta-object but the same container realized as a comonad, so a protocol is a
directed container, hence a small category, hence a polynomial comonad; the cofree comonad on the
tool functor is the interaction trace, and stateful conversational context is a store-comonad
worker. For a single agent this store structure is exactly curried state-threading and buys nothing.
Its value appears only under composition: interleaving two interacting stateful protocols is a
comonad distributive law, equivalently a Zappa--Sz\'ep weld of the underlying directed containers,
which does not exist for free. Once a freeness condition on the two protocols holds, whether a
consistent interleaving exists is governed by a single cohomology class $[\omega]\in H^2$, the weld
existing exactly when $[\omega]=0$; at the smallest re-entrant model the class is one bit
$[\omega]=\varepsilon\in\mathbb{Z}/2$, machine-checked in Lean~4. The note claims no new theorem; its contribution is the identification, and a positioning
against the coalgebraic and effect-handler semantics of protocols that stop one refinement short of
this structure.
\end{abstract}

\maketitle

\begingroup\small
\noindent\fbox{\begin{minipage}{0.95\linewidth}
\emph{Status.} Exploratory work-in-progress note (Kodamai containers series). Every mathematical
fact invoked is proved, and where stated machine-checked, elsewhere; the contribution here is the
synthesis and its honest positioning, not a new theorem. The single result specific to this note's
running model -- the value of the obstruction class at the smallest re-entrant protocol -- is
verified in Lean~4 (\S\ref{sec:compose}, Remark~\ref{rem:lean}). Not yet reviewed.
\end{minipage}}
\endgroup

\bigskip

\section{Introduction}\label{sec:intro}

An agent built on a language model does its work by \emph{calling tools}. It emits a structured
request -- a function name with typed arguments, a database query, a call to another agent -- and
the surrounding runtime, the \emph{harness}, executes that request and returns a structured
response, which the agent folds back into its context before deciding what to do next. The
governing discipline -- which requests are well-formed, which responses each request may receive,
how the responses accumulate into a growing context -- is what recent engineering practice, following
the Model Context Protocol and its relatives, calls a \emph{tool-calling protocol}. Multi-agent
systems compose these: a supervisor agent calls a worker agent as a tool, two agents read and write
a shared store, an agent calls back into the very supervisor that invoked it.

The engineering of such systems is well developed; their \emph{semantics} is not. What kind of
mathematical object is a tool-calling protocol? What object is the harness that runs it? And when
two agents interact, what decides whether their composite behaves consistently, or instead diverges
into the re-entrancy and stale-read defects that plague concurrent stateful systems? These questions
are usually answered operationally, one system at a time, and the correctness of a multi-agent
interleaving is typically discovered the way concurrency bugs always are: at run time, after it
breaks.

\paragraph{The gap.}
There is a natural home for these questions, and it has been standing empty. A container
\cite{gambinokock} is a set of shapes $S$ together with a family of positions $P:S\to\Set$; its
extension is the polynomial functor $\ext{S\ctr P}(X)=\sum_{s:S}(P\,s\to X)$. Reading $S$ as the
request schemas and $P\,s$ as the admissible responses to schema $s$, a single protocol step
\emph{is} a container. This much is not new: a session coalgebra in the sense of Keizer, Basold and
P\'erez \cite{sessioncoalg} is a coalgebra $c:X\to F(X)$ for exactly this endofunctor
$F(X)=\coprod_{a\in A}X^{B_a}$, and they trace the observation to Gambino and Kock. What has not
been done is to carry the identification past the single step. A session coalgebra is a coalgebra
for the protocol's \emph{transition} map; it is not a \emph{comonad} on a protocol object, it carries
no cofree interaction trace, no substitution or tensor algebra of protocols, no notion of stateful
context threaded across steps, and no account of when two protocols compose. The one published line
that treats tool-calling algebraically -- Segura's Lawvere-theory handlers for prompt pipelines
\cite{segura} -- is protocol-agnostic and does not model the protocol object at all. The comonadic,
directed-container layer -- precisely the layer that the theory of directed containers exists to
supply -- is unclaimed.

\paragraph{The observation.}
The missing ingredient is the harness. A tool-calling protocol is not run by a meta-level machine
that stands outside it; the harness is the discipline of threading context from one step to the
next, and that discipline is \emph{structure on the same container}. A container whose extension
carries a comonad structure -- a counit ``answer from the current context'' and a comultiplication
``the context of the context'' -- is exactly a \emph{directed container} in the sense of Ahman and
Uustalu \cite{dircont}, and directed containers are, up to equivalence, small categories, and up to
equivalence polynomial comonads. So the harness upgrades the protocol from a bare container to a
directed container, and the seed equivalence
\begin{equation}\label{eq:chain}
  \DCont\;\simeq\;\Cat\;\simeq\;\Cat^{\sharp}\ \subset\ \Poly
\end{equation}
(directed container $\simeq$ small category $\simeq$ polynomial comonad) does the rest of the lifting
for free. Once the protocol is a comonad, three further structures are immediate rather than
invented: the \emph{cofree comonad} on the tool functor is the interaction trace, the tree of all
runs of the protocol; the store comonad $\Delta S$ models conversational \emph{state}; and the
substitution $\comp$ and Dirichlet tensor $\otimes$ of \cite{spivakref,niuspivak} are the two ways of
composing protocols in sequence and in parallel.

\paragraph{The main point, stated as a dichotomy.}
The state structure is worth its cost only under composition, and the sharp statement is a dichotomy
we import rather than prove. For a single agent, the store worker is exactly hand-threaded state
(Proposition~\ref{prop:single}): a getter and a setter, no more. For two \emph{interacting} stateful
agents -- ones that read and write each other's state -- the composite is neither the free parallel
tensor nor a bare nesting but a \emph{Zappa--Sz\'ep weld}, and a weld is a distributive law saying
how the two contexts interleave. Comonads, unlike the data they carry, do not compose for free
\cite{ahmanuustaluDL}. Whether a consistent interleaving exists is a single cohomology class.

\begin{theorem}[imported; \S\ref{sec:compose}]\label{thm:main}
Let two interacting stateful protocols be presented as directed containers $C,D$. A consistent
interleaving of their harnesses -- a comonad distributive law, equivalently a Zappa--Sz\'ep weld
$C\zs D$ -- exists if and only if a freeness condition \textup{(L)} holds and an obstruction class
\[
  [\omega]\in H^2 \qquad\text{vanishes,}\qquad \text{(G)}\ \Longleftrightarrow\ [\omega]=0 .
\]
At the smallest re-entrant supervisor--worker model the class lives in $H^2\cong\mathbb{Z}/2$ and
equals a single bit,
\[
  [\omega]=\varepsilon\in\mathbb{Z}/2,\qquad
  \varepsilon=\text{``does the callee mutate state the caller depends on?''},
\]
so the weld exists iff $\varepsilon=0$. This value is machine-checked in Lean~4.
\end{theorem}

\noindent The engineering reading is the payoff. A wrong interleaving -- a distributive law that
violates the coherence axioms -- is not an abstract nuisance; it is precisely the state-divergence
and re-entrancy defect, where one agent reads context that another has already invalidated.
Hand-threading the shared state does not escape this: every call site that reads one register and
writes another silently commits to an order of effects, an implicit distributive law, and nothing
checks that the local choices are globally consistent. A program can hand-thread its way into a
globally inconsistent interleaving one locally reasonable line at a time, and that inconsistency is
$[\omega]\ne 0$, found at run time as a bug. The comonadic view does not make the composite exist --
no formalism can, because sometimes it genuinely does not -- but it \emph{concentrates the decision}:
the interleaving is one named object, its soundness is one class, and both can be inspected before
the system runs.

\paragraph{Contribution and honest scope.}
This note proves no new theorem. Its contribution is the identification of the tool-calling protocol
and its harness with the directed-container / polynomial-comonad structure, the resulting reading of
the trace, state, and composition operators, and an honest positioning against the two nearest
precedents \cite{sessioncoalg,segura}. Every mathematical fact it uses is established elsewhere: the
equivalence chain \eqref{eq:chain} is Ahman--Uustalu \cite{dircont}; the distributive-law
$\Leftrightarrow$ Zappa--Sz\'ep correspondence and the $[\omega]\in H^2$ obstruction are from the
directed-container and bicrossed-product literature and our own prior work; the value
$[\omega]=\varepsilon$ at the smallest model is machine-checked (Remark~\ref{rem:lean}). We draw the
formalisation boundary explicitly where it falls (Remark~\ref{rem:lean}), and we flag the one
reduction that is invoked at the directed-container level and not yet restated directly for store
comonads as an open problem (\S\ref{sec:conclusion}).

\paragraph{Outline.}
Section~\ref{sec:prelim} recalls containers, the equivalence chain, the cofree and store comonads,
and the obstruction we import. Section~\ref{sec:protocol} identifies the protocol with a container
and the harness with its comonad, and reads off the interaction trace. Section~\ref{sec:single}
gives the single-agent verdict: state is threaded state, and buys nothing.
Section~\ref{sec:compose} gives the composition verdict: interleaving is a distributive law and its
soundness is $[\omega]=0$, with the Lean boundary. Section~\ref{sec:related} positions the note
against session coalgebras, effect handlers, and process-calculus accounts.
Section~\ref{sec:conclusion} states the open problems.

\section{Preliminaries}\label{sec:prelim}

We recall only what the argument uses; a reader fluent in containers and $\Poly$ can skim to
\S\ref{sec:protocol}. Throughout, $\Set$ is the base and $\y$ denotes the identity container
$1\ctr(\lambda\_.\,1)$, whose extension is the identity functor.

\begin{definition}[Container, extension, morphism \cite{gambinokock}]\label{def:cont}
A \emph{container} $S\ctr P$ consists of a set of \emph{shapes} $S$ and a family of \emph{positions}
$P:S\to\Set$. Its \emph{extension} is the polynomial functor
\[
  \ext{S\ctr P}(X)\;=\;\sum_{s:S}\big(P\,s\to X\big).
\]
A \emph{morphism} $(S\ctr P)\to(T\ctr Q)$ is a pair $(f,\bk f)$ of a forward map $f:S\to T$ on shapes
and a backward, contravariant family $\bk f:\prod_{s:S}\big(Q(f\,s)\to P\,s\big)$ on positions.
Containers and their morphisms form the category $\Cont$; extension is a full and faithful embedding
$\Cont\hookrightarrow[\Set,\Set]$ onto the polynomial functors.
\end{definition}

\begin{definition}[Directed container; the equivalence chain \cite{dircont}]\label{def:dcont}
A \emph{directed container} is a container $S\ctr P$ together with operations making each position
set the arrows out of a shape: a unit $o:\prod_{s}P\,s$, a ``sub-shape'' operation
$s\oplus\_:\prod_s(P\,s\to S)$, and a composition of positions, subject to category-like axioms. The
resulting data is exactly a small category: shapes are objects, $P\,s=\coprod_{t}\Cat(s,t)$ collects
the arrows out of $s$. Passing to extensions, a directed container is precisely a \emph{polynomial
comonad} $(\ext{S\ctr P},\epsilon,\delta)$, with counit $\epsilon:\ext{S\ctr P}\Rightarrow\mathrm{Id}$
and comultiplication $\delta:\ext{S\ctr P}\Rightarrow\ext{S\ctr P}\comp\ext{S\ctr P}$. This gives the
equivalence chain~\eqref{eq:chain}, $\DCont\simeq\Cat\simeq\Cat^{\sharp}$.
\end{definition}

\begin{definition}[Substitution and tensor \cite{spivakref,niuspivak}]\label{def:comp}
For containers $C=S\ctr P$ and $D=T\ctr Q$, \emph{substitution} $C\comp D$ has extension
$\ext C\circ\ext D$ (nest a $D$-structure at each position of a $C$-structure); it is the composition
of the corresponding polynomial functors. The \emph{Dirichlet tensor} $C\otimes D$ multiplies shapes
and positions, $(S\times T)\ctr\big(\lambda(s,t).\,P\,s\times Q\,t\big)$; on extensions it is the
Day convolution for the pointwise product. Substitution models sequential composition; the tensor
models independent parallel composition.
\end{definition}

\begin{definition}[Cofree comonad]\label{def:cofree}
For a container $F$ with extension $\ext F$, the \emph{cofree comonad} $\cofree F$ is the terminal
$F$-coalgebra structure carried by the greatest fixed point $\nu X.\,(A\times\ext F X)$ relative to a
label object $A$; concretely $\cofree F(A)$ is the set of $A$-labelled trees branching by the
positions of $F$. It is the universal comonad mapping to $F$: every run of the one-step functor $F$
extends uniquely to an element of the trace $\cofree F$.
\end{definition}

\begin{definition}[Store comonad and state object \cite{niuspivak}]\label{def:store}
For a set $S$, the \emph{state object} is the container $\Delta S:=S\ctr(\lambda\_.\,S)$, with
extension
\[
  \ext{\Delta S}(X)=\textstyle\sum_{s:S}X^{S}=S\times X^{S}=\Store_S X,
\]
the \emph{store} (costate) comonad: a current state paired with a reader $X^S=(S\to X)$. It is a
directed container -- the codiscrete category on $S$, in which every pair of objects has a unique
arrow. Its counit reads the current state, and its comultiplication duplicates the register. A
\emph{worker of state $S$ from $X$ to $Y$} is a container morphism $w:\Delta S\otimes X\to Y$.
\end{definition}

\noindent Unfolding Definition~\ref{def:cont} for a worker $w:\Delta S\otimes X\to Y$ with
$X=X_0\ctr X_1$ and $Y=Y_0\ctr Y_1$ gives a forward map and a backward family
\begin{equation}\label{eq:fwdbwd}
  \fwd:S\times X_0\to Y_0,
  \qquad
  \bwd:\textstyle\prod_{(s,x_0)}\big(Y_1(\fwd(s,x_0))\to S\times X_1\,x_0\big),
\end{equation}
and, by the universal property of the product, $\bwd$ splits into a \emph{state-writeback}
$\bwd_S:\prod Y_1(\cdots)\to S$ and an ordinary back-map $\bwd_X:\prod Y_1(\cdots)\to X_1\,x_0$. The
forward leg reads state and input and answers; the backward leg writes back a new state and answers
the input. This is the whole of ``a stateful step,'' typed.

\begin{definition}[Zappa--Sz\'ep weld; the obstruction we import \cite{ahmanuustaluDL}]\label{def:zs}
A \emph{Zappa--Sz\'ep weld} $C\zs D$ of two directed containers is a directed-container structure on
a shared underlying object in which each of $C$ and $D$ acts on the other, restricting to $C$ and $D$
on the two factors. Welds correspond to distributive laws
$\lambda:\ext D\comp\ext C\Rightarrow\ext C\comp\ext D$ (equivalently the comonad-distributive dual)
satisfying the coherence axioms; at the level of the underlying categories this is the classical
bicrossed product. A weld exists iff a freeness condition \textup{(L)} holds and a global-closure
condition \textup{(G)} holds, and the obstruction to \textup{(G)} is a single class $[\omega]\in H^2$
with $\text{(G)}\Leftrightarrow[\omega]=0$. We invoke these facts; they are proved in the cited work
and our own.
\end{definition}

\section{A protocol is a container; the harness is its comonad}\label{sec:protocol}

We fix terminology at the level of generality the note needs: a tool-calling protocol is any
discipline of typed requests and typed responses, of which the Model Context Protocol is one
concrete instance. We do not model any particular wire format; we model the shape of the interaction.

\begin{definition}[Tool-calling protocol]\label{def:protocol}
A \emph{tool-calling protocol} is a container $S\ctr P$ in which $S$ is the set of well-formed
\emph{request schemas} an agent may emit (a tool name together with typed arguments, a query, a
delegated sub-task) and, for each $s:S$, $P\,s$ is the set of admissible \emph{responses} the harness
may return to $s$. One \emph{step} of the protocol is the endofunctor $\ext{S\ctr P}(X)=\sum_{s:S}(P\,s\to
X)$: choose a request, and continue as a function of whichever response comes back.
\end{definition}

\begin{remark}[What is conceded]\label{rem:concede}
Definition~\ref{def:protocol}, up to the reading of $S$ and $P$, is not new. The step endofunctor is
the session-coalgebra functor $F(X)=\coprod_{a\in A}X^{B_a}$ of Keizer, Basold and P\'erez
\cite{sessioncoalg}, who study coalgebras $c:X\to F(X)$ for it (they cite Gambino--Kock, and do not
use the word ``container''). We take ``one protocol step is a container'' from them. Everything from
here on -- the comonad, the trace, the store worker, the composition operators, and the obstruction
-- lies outside their development; see \S\ref{sec:related}.
\end{remark}

The step functor alone is only half the object. An agent does not issue one request in a vacuum; the
harness \emph{threads context} -- the growing transcript of prior requests and responses -- from each
step to the next, and it is this threading, not the individual step, that makes a protocol a
protocol. The claim of this section is that the threading is not an external machine but a comonad
structure on the very same container.

\begin{proposition}[The harness is the protocol's comonad]\label{prop:harness}
A tool-calling protocol whose harness threads context coherently -- an initial (empty) context, a
way to read the current context to answer without a further request, and an associative extension of
a context by a step -- is a directed container: the counit $\epsilon$ is ``answer from the current
context,'' the comultiplication $\delta$ is ``expose the context of each continuation,'' and the
category-theoretic axioms of Definition~\ref{def:dcont} are exactly the coherence of context
threading. By the equivalence chain~\eqref{eq:chain} it is therefore a small category and a
polynomial comonad.
\end{proposition}

\begin{proof}
This is the content of Definition~\ref{def:dcont} read in the protocol's vocabulary, not a new
theorem. A coherent threading discipline provides precisely the unit $o$ (the empty context is a
distinguished position at each shape), the sub-shape map $s\oplus\_$ (a response determines the
continuation's schema), and the composition of positions (concatenation of transcripts), subject to
the directed-container axioms, which say that the empty context is a two-sided unit for concatenation
and that concatenation is associative. These are the comonad laws for $(\ext{S\ctr P},\epsilon,\delta)$
under \eqref{eq:chain}.
\end{proof}

\noindent The harness is thus not a meta-object above the protocol; it \emph{is} the protocol, with
the comonad structure that the word ``harness'' names. Two consequences are immediate.

\begin{corollary}[The interaction trace is the cofree comonad]\label{cor:trace}
The interaction trace of a protocol $S\ctr P$ -- the tree of all possible runs, each run a sequence
of request/response pairs -- is the cofree comonad $\cofree{S\ctr P}$ of Definition~\ref{def:cofree}
on the tool functor, with labels drawn from the agent's observable outputs. A single conversation is
one root-to-node branch; the comonad's counit reads the current node, and its comultiplication
re-roots the trace at each node, giving the sub-conversation visible from there.
\end{corollary}

\begin{remark}[Running the agent is an interaction law]\label{rem:interaction}
The comonad reading also locates the agent. An agent's tool-issuing behaviour is a monad of effects,
and executing it against the harness is the pairing of that monad with the protocol comonad by an
\emph{interaction law} in the sense of Katsumata, Rivas and Uustalu \cite{interactionlaws}: the
categorical form of ``running'' the agent against the protocol to produce a value. We do not develop
the agent-monad here -- this note is about the protocol object -- but record that the harness-comonad
identification makes ``run'' a standard construction rather than an operational primitive.
\end{remark}

\begin{corollary}[Protocols compose in two ways]\label{cor:compose}
Two protocols compose by substitution $C\comp D$ -- run $D$ at each response position of $C$, i.e.\ a
tool whose responses are themselves protocols (a supervisor delegating to a sub-protocol) -- and by
Dirichlet tensor $C\otimes D$ -- offer both protocols' requests independently (two non-interacting
tools). These are Definition~\ref{def:comp}. The interesting composite, of two protocols that share
state, is neither, and is the subject of \S\ref{sec:compose}.
\end{corollary}

\section{The single-agent verdict: state is threaded state}\label{sec:single}

Conversational context that the agent may also \emph{mutate} -- a scratchpad, a running plan, a
budget -- is modelled by giving the worker a state object $\Delta S$ (Definition~\ref{def:store}).
Before claiming anything for this, we grant the sceptic the single-agent case in full.

\begin{proposition}[One worker is threaded state, re-spelled]\label{prop:single}
Fix a state set $S$. The elements of the store $\Store_S X=S\times X^S$ are in natural bijection with
state-threading data: a current state together with a map from states to outputs. On the lens
fragment $X=\y$, $Y=A\ctr(\lambda\_.B)$, a worker $\Delta S\to A\ctr(\lambda\_.B)$ is exactly a pair
\[
  \getop:S\to A,\qquad \putop:S\times B\to S,
\]
the ordinary ``read the state to answer; write the state back'' of hand-threaded code.
\end{proposition}

\begin{proof}
The bijection $\Store_S X=S\times X^S$ is the definition of $\ext{\Delta S}$. Specialising
\eqref{eq:fwdbwd} to $X=\y$ (so $X_0=X_1=1$) and $Y=A\ctr(\lambda\_.B)$ gives $\fwd:S\to A$, i.e.\
$\getop$, and $\bwd_S:\prod_s(B\to S)$, i.e.\ $\putop:S\times B\to S$, with $\bwd_X$ trivial.
\end{proof}

\noindent There is no daylight here. A store is a state and a reader; a lens is a getter and a
setter; threading a state through a signature gives a getter and a setter. The comonad's counit is
``read the current state'' ($\getop$) and its comultiplication is ``duplicate the register'' --
operations hand-threaded code already performs. For a system of one stateful agent the store comonad
is a rename: more category theory, no more safety. Honesty requires saying so plainly, and it is the
reason the value must be sought elsewhere. It is not sought in vain, because agentic systems are not
systems of one agent.

\section{The composition verdict: interleaving is a distributive law}\label{sec:compose}

Everything the state structure is worth lives in how two stateful agents combine, and here the store
comonad stops being a rename, because comonads -- unlike the data they carry -- do not compose for
free \cite{ahmanuustaluDL}.

\paragraph{Two ways to stack, only one free.}
Tensoring two state objects multiplies their states on the nose,
\begin{equation}\label{eq:tensor}
  \Delta S\otimes\Delta T=\Delta(S\times T),
\end{equation}
whereas substitution nests a store inside a store,
\begin{equation}\label{eq:nest}
  \Delta S\comp\Delta T=(S\times T^{S})\ctr(\cdots)\;\ne\;\Delta(S\times T),
\end{equation}
recording not a pair of states but a whole branch map $S\to T$; at $|S|=|T|=2$ these are
$4\cdot\y^4$ against $8\cdot\y^4$, an explicit retract, not an equality. Equation~\eqref{eq:tensor}
is the composite of two agents that do \emph{not} interact -- independent registers, no shared reads
or writes -- and it is forced, ``ask both, answer both.'' It is free precisely because it assumes
non-interference.

\paragraph{The interacting composite is neither.}
The agents an orchestration actually composes do interact: a supervisor over a worker that can call
back into the supervisor (re-entrancy), or two agents that read and write a shared ledger, each
acting on data the other produced. There the registers are not independent, so the composite is not
$\Delta S\otimes\Delta T$; and it is not a bare nesting either. It is a Zappa--Sz\'ep weld
$\Delta S\zs\Delta T$ (Definition~\ref{def:zs}), in which each store acts on the other, and a weld is
a distributive law $\lambda$ recording how the two contexts interleave: which agent reads first,
whose write the other sees. Comonads do not compose for free -- without a valid $\lambda$ there is no
composite protocol at all, only two protocols and an unanswered question about the order of their
effects.

\paragraph{Existence is a cohomology class.}
By Theorem~\ref{thm:main}, once (L) holds the weld exists exactly when $[\omega]=0\in H^2$. At the
smallest re-entrant supervisor--worker model the class is one bit
\begin{equation}\label{eq:omega}
  [\omega]=\varepsilon\in\mathbb{Z}/2,\qquad
  \varepsilon=\text{``does the worker mutate shared state the supervisor depends on?''},
\end{equation}
so the weld exists iff $\varepsilon=0$. A \emph{wrong} interleaving -- a $\lambda$ violating the
coherence axioms -- is not an abstract nuisance: it is the state-divergence / re-entrancy defect, the
concrete failure where one agent reads context another has already invalidated.

\begin{remark}[Where the machine check stops]\label{rem:lean}
The class computation~\eqref{eq:omega} is machine-checked in Lean~4, \texttt{sorry}-free and
Mathlib-free (\texttt{Reentrancy.lean}): the file certifies $\varphi(\omega\,\varepsilon)=\varepsilon$
and $\mathsf{InB2}(\omega\,\varepsilon)\Leftrightarrow\varepsilon=\mathsf{false}$, i.e.\ that the
cocycle $\omega(\varepsilon)$ is a coboundary exactly when $\varepsilon=0$. What Lean certifies is the
finite $\mathbb{F}_2$ class computation; the \emph{reduction} of the categorical obstruction to that
cochain complex is proved on paper and is not formalised. The state arithmetic
\eqref{eq:tensor}--\eqref{eq:nest} is likewise Lean-checked (\texttt{WorkersRetract.lean}, the retract
$r\circ\sigma=\mathrm{id}$). We flag the boundary rather than blur it. The general reduction of
store-comonad interleaving to a distributive law is invoked here at the directed-container level --
$\Delta S$ is a directed container -- where it is available; a restatement directly for store comonads
is not needed for this note and is recorded as an open problem in \S\ref{sec:conclusion}.
\end{remark}

\paragraph{What threading does, and does not, do.}
Hand-threading two shared states does not escape the distributive law. Every call site that reads one
register and writes another silently commits to an order of effects -- an implicit $\lambda$ -- and
nothing checks that the per-site choices are globally consistent, which is exactly what the coherence
axioms demand. A program can hand-thread its way into a globally inconsistent interleaving one locally
reasonable line at a time; that inconsistency is $[\omega]\ne0$, discovered at run time as a
re-entrancy bug. The store-comonad worker does not make the composite exist automatically -- no
formalism can, because sometimes it genuinely does not -- but it \emph{concentrates} the decision: the
interleaving is one named object $\lambda$, its validity is one class $[\omega]$, and both can be
inspected before the system runs, in one place, instead of being reconstructed from scattered call
sites after it breaks.

\paragraph{The ledger.}
Table~\ref{tab:ledger} states the verdict as plainly as the question deserves. The reformulation buys
nothing for a single agent, and it does not make composition automatic; what it buys is a change in
\emph{where} and \emph{when} the one real decision of a composite -- the interleaving of shared
effects -- is made and checked.

\begin{table}[t]
\small
\begin{tabular}{@{}p{0.24\linewidth}p{0.30\linewidth}p{0.36\linewidth}@{}}
\toprule
 & \textbf{single agent} & \textbf{interacting composite} \\
\midrule
store-comonad worker & $\Store_S X=S\times X^S$ & Zappa--Sz\'ep weld; needs a law $\lambda$ \\
hand-threaded state & getter / setter & implicit $\lambda$ per call site \\
\textbf{what the comonad adds} & \emph{nothing} (a rename) & the interleaving as \emph{one} named
$\lambda$; its validity as \emph{one} class $[\omega]$, checkable before the system runs \\
\bottomrule
\end{tabular}
\medskip
\caption{What the store-comonad reformulation buys, by system size.}\label{tab:ledger}
\end{table}

\section{Positioning}\label{sec:related}

The note sits one refinement past the two nearest semantic accounts of protocols, and beside a
process-calculus line.

\paragraph{Session coalgebras.}
Keizer, Basold and P\'erez \cite{sessioncoalg} give a coalgebraic semantics of session types over
exactly the container endofunctor $F(X)=\coprod_{a\in A}X^{B_a}$, which we adopt for the single step
(Remark~\ref{rem:concede}). Their object is a coalgebra $c:X\to F(X)$ for protocol \emph{transitions}
-- a single one-step map, with state folded into the carrier $X$ -- and their composition is the
$\pi$-calculus parallel product $P\mid Q$. They introduce no comonad, no cofree interaction trace, no
substitution or tensor algebra of protocols, no store-comonad state, and no obstruction to
composition. The comonadic / directed-container layer, the cofree trace, the $\comp$/$\otimes$
algebra, the store worker, and the $[\omega]$ obstruction are outside their development, and are what
this note contributes.

\paragraph{Effect handlers for prompt pipelines.}
Segura \cite{segura} models bounded LLM prompt pipelines with Lawvere theories and algebraic-effect
handlers (retrieval, state, exceptions, typed tool invocation, with live/log/replay interpretations).
This is the closest ``algebra of tool-calling,'' but it is protocol-agnostic: it axiomatises the
effects a pipeline may perform, not the protocol object through which requests and responses flow,
and it uses no polynomial-functor or container structure. We know this work at abstract level only
(the full text is paywalled and we have not read it in depth); we position against it rather than
build on it, and cite it as the effect-handler precedent that does not model the protocol object.

\paragraph{Process-calculus and other accounts.}
Schlapbach's process-calculus treatment of protocol bisimulation is a natural target for a
directed-container morphism reformulation, which we do not pursue here. A withdrawn attempt at a
``monadic'' context-engineering semantics (Zhang--Wang, arXiv:2512.22431, withdrawn for issues in its
categorical foundations) indicates that the naive monadic route does not reach the structure; the
comonadic route of this note is, we believe, the correct one, supplied for free by the equivalence
chain.

\section{Conclusion and open problems}\label{sec:conclusion}

A tool-calling protocol is a container, its harness is that container's comonad, its interaction
trace is the cofree comonad, its stateful context is a store worker, and the correctness of a
multi-agent interleaving is the vanishing of one cohomology class. None of these is a new theorem;
together they say that the object applied category theory already studies -- the directed container,
equivalently the polynomial comonad -- is the semantics of agent tool-calling, and that the
multi-agent correctness question is not a testing problem but a computable obstruction. For an
orchestration of many interacting agents, which is what an agentic system is, that obstruction is the
compositional-correctness guarantee worth naming.

Three problems remain open.

\begin{enumerate}[leftmargin=1.6em]
\item \emph{A store-comonad restatement of the obstruction.} Theorem~\ref{thm:main} is invoked at the
directed-container level, where $\Delta S$ lives. Restating the reduction ``store-comonad interleaving
$\Leftrightarrow$ distributive law $\Leftrightarrow[\omega]$'' directly in the store-comonad
vocabulary, and formalising the categorical-to-cochain reduction that Remark~\ref{rem:lean} leaves on
paper, is a concrete next step.
\item \emph{Beyond the smallest model.} The bit $[\omega]=\varepsilon\in\mathbb{Z}/2$ is the
supervisor--worker case. The class for a weld with a larger freeness datum (three or more interacting
agents, or non-abelian shared state) is the natural generalisation; the freeness condition (L) is what
must first be understood in the many-agent case.
\item \emph{Protocol morphisms.} Reformulating protocol refinement and bisimulation -- including the
process-calculus accounts -- as directed-container morphisms would connect this semantics to the
existing operational theory and is the route to comparing two harnesses for the same protocol.
\end{enumerate}

\subsection*{Acknowledgements}
For Robin and Neil. This note answers a question Robin posed about whether the store-comonad worker
earns its keep, recast for an external audience; the internal house-doc version is companion material.

\begin{thebibliography}{9}

\bibitem{dircont}
D.~Ahman and T.~Uustalu.
\emph{Directed containers as categories.}
Electron.\ Proc.\ Theor.\ Comput.\ Sci.\ (EPTCS) 207 (2016), 89--98. arXiv:1604.01187.

\bibitem{ahmanuustaluDL}
D.~Ahman and T.~Uustalu.
\emph{Distributive laws of directed containers.}
Progress in Informatics 10 (2013), 3--18.

\bibitem{gambinokock}
N.~Gambino and J.~Kock.
\emph{Polynomial functors and polynomial monads.}
Math.\ Proc.\ Cambridge Philos.\ Soc.\ 154 (2013), no.~1, 153--192. arXiv:0906.4931.

\bibitem{sessioncoalg}
A.~C.~Keizer, H.~Basold, and J.~A.~P\'erez.
\emph{Session coalgebras: a coalgebraic view on session types and communication protocols.}
ACM Trans.\ Program.\ Lang.\ Syst.\ (TOPLAS), 2021. arXiv:2011.05712.

\bibitem{niuspivak}
N.~Niu and D.~I.~Spivak.
\emph{Polynomial functors: a mathematical theory of interaction.}
arXiv:2312.00990, 2023.

\bibitem{spivakref}
D.~I.~Spivak.
\emph{A reference for categorical structures on \textup{Poly}.}
arXiv:2202.00534, 2022.

\bibitem{segura}
J.~J.~Segura.
\emph{Algebraic effects for bounded prompt pipelines: Lawvere theories, handlers, and outcome
traces.}
Computing 108 (2026). DOI 10.1007/s00607-026-01699-w. \textup{(Known to the author at abstract level
only; full text paywalled.)}

\bibitem{interactionlaws}
S.~Katsumata, E.~Rivas, and T.~Uustalu.
\emph{Interaction laws of monads and comonads.}
LICS 2020. arXiv:1912.13477.

\end{thebibliography}

\end{document}
